\subsection{§III.6.3, Corollaire 1 --- $\mathfrak{a}^n = \mathfrak{a}$, et ce que
coûte une promesse en commentaire}
\label{sec:iii6-cor1}

\textbf{Énoncé (Bourbaki, E III.49).} Si $\mathfrak{a}$ est un cardinal infini et
$n \ge 1$ un entier, $\mathfrak{a}^n = \mathfrak{a}$. La démonstration est une
récurrence sur $n$ à partir de $1$, dont l'hérédité repose sur le Théorème 2.

\textbf{Ce que le dépôt contenait déjà.} Toute la récurrence :
$\mathcal{F}(n;A)$ borné inférieurement puis supérieurement, Cantor--Bernstein
pour conclure, base $R\{1\}$ close, hérédité $R\{n\} \Rightarrow R\{n+1\}$.
Le travail mathématique était fait. Mais \texttt{a\_puissance\_n\_egale\_a}
concluait sous \textbf{trois} hypothèses --- \texttt{predecesseur\_fini\_universel},
$\mathrm{Eq}(A\times A, A)$, et $a_0 \in A$ --- et sa docstring annonçait la
décharge de la deuxième « via \texttt{hessenberg\_0hyp} ».

\textbf{Cette fonction n'existait pas.} Un \texttt{grep} sur le dépôt entier en
donne une seule occurrence : celle de la docstring qui la promet. C'est le même
défaut que la surdéclaration débusquée dans le dossier Goldbach, en plus
discret : \emph{un nom de fonction dans un commentaire n'est pas une fonction},
et aucun test ne peut attraper cela --- le noyau garantit la soundness, jamais
la fidélité d'une prose.

\textbf{Les trois décharges.} Chacune a demandé un objet distinct, et aucune
n'était où on l'attendait.

\emph{(i) Le carré.} Le Théorème 2 conclut sur les \textbf{cardinaux}
($\mathrm{Card}\,A \cdot \mathrm{Card}\,A = \mathrm{Card}\,A$) ; le corollaire
réclame les \textbf{ensembles} ($\mathrm{Eq}(A\times A, A)$). Le pont tient en
trois pas : \texttt{produit\_cardinal\_bien\_defini} nourri de deux réflexivités
donne $\mathrm{Card}(A\times A) = \mathrm{Card}\,A \cdot \mathrm{Card}\,A$
(clos, $1{,}2$~s) ; Hessenberg donne le membre droit ; la Proposition 1 dans son
sens réciproque conclut. Le test vérifie l'identité de formules par
\texttt{==} et non la ressemblance --- une ressemblance ne déchargerait rien.

\emph{(ii) Le point $a_0$.} Bourbaki ne mentionne pas cette hypothèse, et il a
raison : elle découle de « $A$ infini ». Encore fallait-il la dériver. La route
est entièrement par contraposition et tient en trois théorèmes \textbf{clos} :
$\emptyset$ est fini, donc un infini n'est pas $\emptyset$, donc il a un élément.
Le point élégant est que $\mathrm{est\_infini}$ est \emph{défini} comme
$\neg\mathrm{Fini}$ (E III.6.1, Déf.~1) : la contraposée tombe exactement sur la
forme voulue, sans pont ni réécriture.

\emph{(iii) Le résidu.} \texttt{predecesseur\_fini\_universel} était --- une fois
de plus --- \textbf{déjà prouvé} par le dépôt et simplement non branché.

\textbf{Une condition de licéité rendue vérifiable.} $a_0 \in A$ disparaît par
$\exists$-élimination, qui n'est licite que si $a_0$ n'est libre dans
\emph{aucune autre} hypothèse au moment où on l'applique. Plutôt que de le
supposer, un étage intermédiaire (\texttt{puissance\_n\_sous\_hessenberg}) porte
cette assertion, et son test la vérifie : si elle tombait, le théorème final
serait invalide sans que rien d'autre ne le signale.

\textbf{Résultat.} $\vdash$ $\mathrm{est\_infini}(\mathrm{Card}\,A) \Rightarrow
(\forall n)\big((n \text{ entier} \wedge n \ge 1) \Rightarrow
\mathrm{Eq}(\mathcal{F}(n;A), A)\big)$ sous la seule hypothèse
\texttt{enonce\_hessenberg(A)} --- le Théorème 2 lui-même. Mesuré : $66{,}4$~s,
une hypothèse, non vacuous, \texttt{theorie\_ensembles()} $= 22$.

Le Théorème 2 est \emph{porté} plutôt que dérivé, par application de la règle de
performance du projet : sa preuve coûte $5172{,}8$~s, et un grand lemme clos ne
se repaie pas à chaque brique intermédiaire --- il devient une hypothèse
honnête, déchargée une fois au dernier étage. \textbf{Statut : FAIT}, au même
titre que le théorème dont il dépend.
